\documentclass[11pt]{article}

\usepackage{amsmath}
\usepackage{amssymb}
\usepackage{amsthm}
\usepackage{booktabs}
\usepackage[hidelinks]{hyperref}

\newtheorem{definition}{Definition}
\newtheorem{proposition}{Proposition}

\newcommand{\cosim}{\sigma}
\newcommand{\thresh}{\tau}

\title{The Mathematics of the Pando Connectivity Engine}
\author{Pando}
\date{July 9, 2026}

\begin{document}

\maketitle

\begin{abstract}
Pando is a word game: connect a start word to an end word through a chain of
intermediate words, where every adjacent pair must be semantically related, and
the shortest verified chain wins (tops the leader board). This note formalizes the mathematics behind
the verification engine: the embedding space and similarity measure, the link
verification predicate, the empirical calibration of the acceptance threshold,
the construction of the candidate word pool and its connectivity graph, and the
breadth-first benchmark solver that publishes a deterministic shortest route
for each puzzle.
\end{abstract}

\section{Introduction}

A Pando puzzle is a pair of words $(s, e)$, with a starting word $s$ 
and a target, ending word $e$. A player answers with a chain of
connector words $c_1, \dots, c_k$, claiming that the sequence
$s \to c_1 \to \dots \to c_k \to e$ moves between semantically connected words
at every step. The engine verifies each adjacent link independently: there are
no pre-programmed routes, no requirement that a link move ``toward'' the
target, and no bound on chain length. Creative detours are legal; the score is
simply $k$, the number of connectors, and fewer is better. The engine also
solves each puzzle itself and publishes its shortest verified route as a
benchmark that players may attempt to beat. 

\section{Embedding space and similarity}

Words are represented by the pretrained \emph{word2vec-google-news-300}
embeddings \cite{mikolov2013}, which assign each vocabulary entry $w$ a vector
$v_w \in \mathbb{R}^{300}$. The engine retains the $500{,}000$ most frequent
entries of the released model; beyond that point the vocabulary is dominated by
rare phrase tokens. All comparisons use unit-normalized vectors
\[
  \hat{v}_w = \frac{v_w}{\lVert v_w \rVert_2},
\]
so that the similarity of two words is the cosine of the angle between their
embeddings, computed as a single dot product:
\begin{definition}[Connectivity]
For words $a, b$ with embeddings in the model,
\[
  \cosim(a, b) = \hat{v}_a \cdot \hat{v}_b \in [-1, 1].
\]
\end{definition}
Larger $\cosim$ means the words occur in more similar contexts in the training
corpus (roughly $100$ billion words of news text), which is the engine's
operational notion of ``word connectivity.''

\section{The link predicate}

A single link $(a, b)$ is \emph{verified} --- signed off with its similarity
value --- exactly when it clears a calibrated threshold and is not a trivial
variant pair, with one deliberate exception for multi-word entities.

\begin{definition}[Verified link]
\label{def:link}
Let $\thresh$ be the acceptance threshold. The link $(a,b)$ verifies iff
\[
  \neg\,\mathrm{var}(a, b)
  \;\wedge\;
  \bigl( \cosim(a, b) \ge \thresh \;\vee\; \mathrm{colloc}(a, b) \bigr),
\]
where $\mathrm{var}$ and $\mathrm{colloc}$ are defined below. A link whose
words are not both in the vocabulary is rejected outright.
\end{definition}

\paragraph{Inflection guard.}
Without a guard, chains like \emph{rain} $\to$ \emph{rains} $\to$
\emph{raining} would verify trivially, since inflected forms of a word are
nearly parallel in embedding space. Let $S(w)$ be the set of plausible stems of
$w$ produced by a deterministic suffix-stripping rule: casefold $w$, then for
each suffix in the family
$\{\text{-s}, \text{-es}, \text{-ed}, \text{-ing}, \text{-er}, \text{-est},
\text{-ies}, \text{-ied}, \text{-ly}, \dots\}$
whose removal leaves at least three characters, add the stripped stem, a
doubled-consonant repair (\emph{running} $\to$ \emph{runn} $\to$ \emph{run}),
and a \mbox{-ies $\to$ -y} repair (\emph{cities} $\to$ \emph{city}). Then
\[
  \mathrm{var}(a, b) \iff
  a \equiv b \ \text{(casefolded)} \ \vee\ S(a) \cap S(b) \neq \emptyset .
\]
Variant pairs are rejected regardless of their cosine.

\paragraph{Collocation rule.}
Pop-culture pairings often live in the model as multi-word \emph{entity}
tokens rather than in the single lowercase tokens: the plain tokens give
$\cosim(\text{star}, \text{wars}) \approx 0.02$, far below any sensible
threshold, yet \texttt{Star\_Wars} exists as its own vocabulary entry. The
predicate
\[
  \mathrm{colloc}(a, b) \iff
  \text{some cased joining of } a, b
  \text{ (either order) is a vocabulary entry}
\]
therefore verifies such pairs as an alternative path, independent of their
token-level cosine.

\section{Threshold calibration}

The threshold is measured instead of chosen a priori. Fix a must-pass set $P$
of word pairs that any reasonable player would accept as connected
(\emph{rock--band}, \emph{apple--pie}, \emph{winter--snow}, \dots) and a
must-fail set $F$ of pairs that no one would (\emph{carrot--algebra},
\emph{tuba--vinegar}, \dots). All cosines are measured against the actual
model, and $\thresh$ must separate the two sets with a margin $m$ on both
sides:
\[
  \min_{(a,b) \in P} \cosim(a,b) \;\ge\; \thresh + m,
  \qquad
  \max_{(a,b) \in F} \cosim(a,b) \;\le\; \thresh - m .
\]
With $m = 0.015$ the calibrated value is
\[
  \boxed{\thresh = 0.19}.
\]
Pairs that verify only through the collocation rule are held in a separate
calibration category exempt from the cosine inequality --- \emph{star--wars}
is a must-pass pair, but at cosine $0.02$ it would make the two sets
inseparable if forced through the threshold. Representative measurements:

\begin{center}
\begin{tabular}{llc}
\toprule
Pair & Role & $\cosim$ \\
\midrule
winter -- snow    & must pass                & $0.490$ \\
money -- bank     & must pass                & $0.261$ \\
rain -- cloud     & must pass (lowest)       & $0.210$ \\
tuba -- vinegar   & must fail (highest)      & $0.160$ \\
carrot -- algebra & must fail                & $0.103$ \\
star -- wars      & collocation (exempt)     & $0.021$ \\
\bottomrule
\end{tabular}
\end{center}

\section{The candidate pool and connectivity graph}

Connector words must come from a canonical pool $P^\ast$ shared by players and
the solver, so the published benchmark is achievable with exactly the moves
players are allowed. The pool is built in two passes.

\paragraph{Frequency filter.}
Take the vocabulary in descending frequency order and keep entries matching
\verb|^[a-z]{2,}$| (clean lowercase alphabetic words) that are not function
words, until $50{,}000$ words are collected. Call this $P_0$.

\paragraph{Hub suppression.}
Generic words (\emph{thing}, \emph{good}, \dots) are close to nearly
everything and would make every puzzle trivially solvable through them. Define
the above-threshold degree of each word,
\[
  d(w) = \bigl|\{\, u \in P_0 \setminus \{w\} : \cosim(w, u) \ge \thresh \,\}\bigr|,
\]
computed exactly by blockwise products of the row-normalized pool matrix
$M \in \mathbb{R}^{|P_0| \times 300}$ against $M^{\mathsf T}$. Words whose
degree exceeds the $99$th percentile of the degree distribution are dropped:
\[
  P^\ast = \{\, w \in P_0 : d(w) \le \operatorname{percentile}_{99}(d) \,\}.
\]
A percentile cap is used rather than an absolute one because the degree
distribution shifts wholesale with $\thresh$; at a permissive threshold nearly
every word exceeds any fixed count.

\paragraph{The graph.}
The pool induces an undirected connectivity graph $G = (P^\ast, E)$ with
\[
  E = \bigl\{\, \{u, w\} : \cosim(u, w) \ge \thresh
  \ \wedge\ \neg\,\mathrm{var}(u, w) \,\bigr\}.
\]
Empirically, at $\thresh = 0.19$ this graph is dense: the median word has
approximately $2{,}520$ above-threshold neighbors. A direct consequence is
that most valid puzzles admit a shortest route with a \emph{single}
intermediate word --- a fact with design implications for a shortest-wins
leaderboard, since ties at $k = 1$ will be common and a secondary criterion
(Section~\ref{sec:solver}) is needed to order them.

\section{Sequence verification}

\begin{definition}[Verified sequence]
For a puzzle $(s, e)$ and connectors $c_1, \dots, c_k$, write the chain
$c_0 = s$, $c_{k+1} = e$. The sequence verifies iff
\begin{enumerate}
  \item every adjacent link $(c_i, c_{i+1})$, $0 \le i \le k$, verifies per
        Definition~\ref{def:link};
  \item every connector lies in the canonical pool: $c_i \in P^\ast$ for
        $1 \le i \le k$ (endpoints need only be in the vocabulary);
  \item the chain words are pairwise distinct (casefolded).
\end{enumerate}
The sequence's score is its length $k$; fewer connectors is better.
\end{definition}
Each link in the result carries its cosine and, when rejected, a machine-readable
reason (\emph{below-threshold}, \emph{not-in-pool}, \emph{not-in-vocabulary},
\emph{inflection-variant}, \emph{repeated-word}), so a failed submission
identifies exactly which link broke.

\section{The benchmark solver}
\label{sec:solver}

The published benchmark route is the shortest verified route through $G$,
found by level-synchronous breadth-first search from $s$ with a hop limit
$H = 8$. Each BFS level is one batched matrix operation: for a frontier of
pool indices $F_t$, the similarities of all frontier words against the entire
pool are the rows of $M_{F_t} M^{\mathsf T}$, and the next frontier is every
unvisited, unblocked word whose similarity to some frontier word clears
$\thresh$ (variant pairs excluded). Every shortest-depth parent of each
discovered word is recorded, so all shortest paths can be enumerated by
backtracking once any frontier word links to $e$.

\paragraph{Cosine edges only.}
The solver deliberately searches only threshold edges, never collocation
edges. A player who spots a collocation shortcut (\emph{star} $\to$
\emph{wars}) can therefore legitimately beat the published benchmark --- this
asymmetry is intentional and funds the game's bounty mechanic.

\paragraph{Determinism.}
Among routes of minimal length $k^\ast$, the solver ranks by total link cosine
\[
  T(c_1, \dots, c_{k^\ast})
  = \sum_{i=0}^{k^\ast} \cosim(c_i, c_{i+1}),
\]
descending, breaking remaining ties by lexicographic order of the connector
tuple. The published route is thus identical across runs and machines. The
total-cosine criterion doubles as the natural secondary leaderboard metric
when many players tie at the minimal length.

\section{Solution trees and puzzle validity}

The top-$r$ shortest routes (all of length $k^\ast$) are merged into a
\emph{solution tree} rooted at $s$: routes sharing a prefix share the
corresponding path of nodes, every edge carries its signed-off link, and every
root-to-leaf path is one of the routes, ending at $e$ and re-verifying as a
sequence. The tree is the engine's certificate that the benchmark and its
near-alternates are all legal play.

Finally, a pair $(s, e)$ is a \emph{publishable puzzle} iff
\begin{enumerate}
  \item both $s$ and $e$ are in the vocabulary;
  \item the pair is not trivially connected: $\cosim(s, e) < \thresh$, the
        pair is not a collocation, and $\neg\,\mathrm{var}(s, e)$;
  \item at least one verified route exists within $H$ hops.
\end{enumerate}
Condition 2 rules out puzzles solvable with zero connectors; condition 3 rules
out unsolvable ones. Each seed puzzle shipped with the engine passes this
gate against the real model.

\begin{thebibliography}{1}
\bibitem{mikolov2013}
T.~Mikolov, I.~Sutskever, K.~Chen, G.~Corrado, and J.~Dean.
\newblock Distributed representations of words and phrases and their
  compositionality.
\newblock In \emph{Advances in Neural Information Processing Systems~26},
  2013.
\end{thebibliography}

\end{document}
